\documentclass{article}
\usepackage{qilin}
\hfuzz=1000pt 
\usepackage{amssymb}
\hbadness=99999 % we're bad students
\hfuzz=100pt % wide bois begone

\usepackage{mathtools}
\usepackage{arydshln}

% \newcommand{\dim}[1]{\mathrm{dim}\,#1}

\title{Adjoint Properties}
\author{QiLin Xue}
\lhead{MAT185}
\rhead{QiLin Xue}

\renewcommand{\qedsymbol}{\includegraphics[height=2ex]{medici.png}}

\begin{document}
    \maketitle
    \section{Definition}
    The adjoint, or adjugate of $\bff{A}$, is the transpose of the cofactor matrix $\bff{C}=[c_{ij}]$:
    \begin{equation*}
        \adj \bff{A} = \bff{C}^T
    \end{equation*}
    where:
    \begin{equation*}
        c_{ij}(\bff{A}) = (-1)^{i+j}\det \bff{M}_{ij}(\bff{A})
    \end{equation*}
    where $\bff{M}$ is the $(i,j)$-minor.
    \section{Best Property}
    The most important property is Theorem VIII in Medici:
    \begin{equation*}
        \bff{A}\adj \bff{A} = (\adj \bff{A})\bff{A} = (\det \bff{A})\bff{I}
    \end{equation*}
    which will turn out to be very useful.
    \section{Other Properties}
    For any $n\times n$ matrix:
    \begin{enumerate}
        \item $\adj \bff{0} = \bff{0}$ and $\adj \bff{I} = \bff{I}$
        \begin{proof}
            To prove the first part, apply the definition and note that all minors are the zero matrix and the determinant of the zero matrix is zero. For the second part, apply Theorem VIII.
        \end{proof}
        \item $\adj \lambda \bff{A} = \lambda^{n-1} \adj \bff{A}$
        \begin{proof}
            Using the definition, note that each entry of the cofactor is:
            \begin{equation*}
                c_{ij}(\lambda \bff{A}) = (-1)^{i+j} \det \bff{M}_{ij}(\lambda \bff{A}) = (-1)^{i+j} \lambda^{n-1}\det \bff{M}_{ij}(\bff{A})
            \end{equation*}
            Since each minor is an $(n-1)\times (n-1)$ matrix. We can factor this out such that:
            \begin{equation*}
                \bff{C}(\lambda \bff{A}) = \lambda \bff{C}(\bff{A}) \implies \adj(\lambda \bff{A}) = \lambda^{n-1}\adj \bff{A}
            \end{equation*}
        \end{proof}
        \item $\adj(\bff{A}^T) = \adj(\bff{A})^T$
        \begin{proof}
            From Theorem VIII, taking the transpose of both sides:
            \begin{equation*}
                \bff{A}^T\adj(\bff{A})^T = \det (\bff{A})\bff{I}
            \end{equation*}
            Now, apply Theorem VIII to $\bff{A}^T$ to get:
            \begin{equation*}
                \bff{A}^T\adj(\bff{A}^T) = \det(\bff{A}^T)\bff{I}
            \end{equation*}
            Since we have:
            \begin{equation*}
                \det(\bff{A}^T)=\det(\bff{A}),
            \end{equation*}
            it means that:
            \begin{equation*}
                \bff{A}^T\adj(\bff{A})^T = \bff{A}^T\adj(\bff{A}^T)
            \end{equation*}
            If $\bff{A}$ is not the zero matrix, then this directly proves it. If $\bff{A}$ is the zero matrix, then from property $1$, we get: $\adj(\bff{A}^T)=\bff{0}$ and $\adj(\bff{A})^T = \bff{0}^T=\bff{0}$.
        \end{proof}
        \item $\det(\adj \bff{A})= \det(\bff{A})^{n-1}$
        \begin{proof}
            Taking the determinant of both sides in Theorem VIII:
            \begin{align*}
                \det(\bff{A}\adj\bff{A})&=\det((\det \bff{A})\bff{I}) \\ 
                \det(\bff{A})\det{\adj\bff{A}}&=\det(\bff{A})^{n-1} \\ 
                \det(\adj \bff{A}) &= \det(\bff{A})^{n-1}
            \end{align*}
        \end{proof}
        \item $\adj(\bff{A}\bff{B})=\adj(\bff{B})\adj(\bff{A})$ where $\bff{B}$ is a $n\times n$ matrix.
        \begin{proof}
            Assume both matrices are invertible. Applying Theorem VIII to expand $\adj\bff{B}\adj\bff{A}$, we get:
            \begin{align*}
                \adj\bff{B}\adj\bff{A} &= (\det \bff{B})\bff{B}^{-1}(\det \bff{A})\bff{A}^{-1} \\ 
                &= (\det \bff{AB})(AB)^{-1} \\ 
                &= \adj(\bff{AB})
            \end{align*}
            If the matrices are non-invertible, then... idk you're fucked I guess. (but trust me, this still works)
        \end{proof}
        \item Using the above properties, we can show that if $\bff A$ is:
        \begin{itemize}
            \item upper triangular
            \item lower triangular
            \item diagonal
            \item orthogonal
            \item symmetric
        \end{itemize}
        then $\adj \bff{A}$ is as well.
        \item If $\rank \bff{A} \le n-1$, then $\rank \adj \bff{A} \le 1$
        \begin{proof}
            Let $\rank \bff{A} = n-1$. Then:
            \begin{equation*}
                \nullity A = n - \rank A = 1
            \end{equation*}
            Since $\bff{A}$ is singular, we have $\det \bff{A}=0$. From Theorem VIII:
            \begin{equation*}
                \bff{A}\adj\bff{A} = 0 \implies \col \adj \bff{A} \in \nulll A
            \end{equation*}
            which implies
            \begin{equation*}
                \nullity \bff{A} \ge \rank \adj \bff{A} \implies 1 \ge \rank \adj \bff{A}
            \end{equation*}
            In fact, the rank will always be $1$, since there will be some non-zero minor. However, I don't know how to prove this last part.
            \vspace{2mm}

            If $\rank \bff{A} \le n-2$, then $\rank \adj \bff{A} = 0.$ This means that the minimum amount of rows that we need to remove such that all rows are linearly independent is $2$. All minors will have $n-1$ rows, so all minors are consisted of linearly dependent rows and as a result, the determinant of the minors will all be zero and thus $\adj \bff{A} = \bff{0}$. 
        \end{proof}
        \item $\adj \adj \bff{A} = \det(\bff{A})^{n-2}\bff{A}$
        \begin{proof}
            Applying Theorem VIII on $\adj\bff{A}$ gives:
            \begin{equation*}
                \adj \bff{A} \adj\adj\bff{A} = (\det \adj \bff{A})\bff{I}
            \end{equation*}
            From property $4$, we have:
            \begin{equation*}
                \adj\bff{A}\adj\adj\bff{A} = \det(\bff{A})^{n-1}
            \end{equation*}
            Multiplying both sides by $\bff{A}$ and applying Theorem VIII again gives:
            \begin{align*}
                (\bff{A} \adj \bff{A})\adj\adj \bff{A} &= \bff{A}\det(\bff{A})^{n-1} \\ 
                \det(\bff{A})\adj\adj\bff{A} &= \det(\bff{A})^{n-1} \\
                \adj\adj\bff{A} &= \bff{A}\det(\bff{A})^{n-2} 
            \end{align*}
        \end{proof}
    \end{enumerate}
\end{document}
